% Options for packages loaded elsewhere
% Options for packages loaded elsewhere
\PassOptionsToPackage{unicode}{hyperref}
\PassOptionsToPackage{hyphens}{url}
\PassOptionsToPackage{dvipsnames,svgnames,x11names}{xcolor}
%
\documentclass[
  russian,
]{article}
\usepackage{xcolor}
\usepackage[top=20mm,left=20mm,right=20mm,bottom=20mm]{geometry}
\usepackage{amsmath,amssymb}
\setcounter{secnumdepth}{5}
\usepackage{iftex}
\ifPDFTeX
  \usepackage[T1]{fontenc}
  \usepackage[utf8]{inputenc}
  \usepackage{textcomp} % provide euro and other symbols
\else % if luatex or xetex
  \usepackage{unicode-math} % this also loads fontspec
  \defaultfontfeatures{Scale=MatchLowercase}
  \defaultfontfeatures[\rmfamily]{Ligatures=TeX,Scale=1}
\fi
\usepackage{lmodern}
\ifPDFTeX\else
  % xetex/luatex font selection
  \setmainfont[]{Times New Roman}
  \setsansfont[]{Arial}
  \setmonofont[]{Courier New}
\fi
% Use upquote if available, for straight quotes in verbatim environments
\IfFileExists{upquote.sty}{\usepackage{upquote}}{}
\IfFileExists{microtype.sty}{% use microtype if available
  \usepackage[]{microtype}
  \UseMicrotypeSet[protrusion]{basicmath} % disable protrusion for tt fonts
}{}
\makeatletter
\@ifundefined{KOMAClassName}{% if non-KOMA class
  \IfFileExists{parskip.sty}{%
    \usepackage{parskip}
  }{% else
    \setlength{\parindent}{0pt}
    \setlength{\parskip}{6pt plus 2pt minus 1pt}}
}{% if KOMA class
  \KOMAoptions{parskip=half}}
\makeatother
% Make \paragraph and \subparagraph free-standing
\makeatletter
\ifx\paragraph\undefined\else
  \let\oldparagraph\paragraph
  \renewcommand{\paragraph}{
    \@ifstar
      \xxxParagraphStar
      \xxxParagraphNoStar
  }
  \newcommand{\xxxParagraphStar}[1]{\oldparagraph*{#1}\mbox{}}
  \newcommand{\xxxParagraphNoStar}[1]{\oldparagraph{#1}\mbox{}}
\fi
\ifx\subparagraph\undefined\else
  \let\oldsubparagraph\subparagraph
  \renewcommand{\subparagraph}{
    \@ifstar
      \xxxSubParagraphStar
      \xxxSubParagraphNoStar
  }
  \newcommand{\xxxSubParagraphStar}[1]{\oldsubparagraph*{#1}\mbox{}}
  \newcommand{\xxxSubParagraphNoStar}[1]{\oldsubparagraph{#1}\mbox{}}
\fi
\makeatother


\usepackage{longtable,booktabs,array}
\newcounter{none} % for unnumbered tables
\usepackage{calc} % for calculating minipage widths
% Correct order of tables after \paragraph or \subparagraph
\usepackage{etoolbox}
\makeatletter
\patchcmd\longtable{\par}{\if@noskipsec\mbox{}\fi\par}{}{}
\makeatother
% Allow footnotes in longtable head/foot
\IfFileExists{footnotehyper.sty}{\usepackage{footnotehyper}}{\usepackage{footnote}}
\makesavenoteenv{longtable}
\usepackage{graphicx}
\makeatletter
\newsavebox\pandoc@box
\newcommand*\pandocbounded[1]{% scales image to fit in text height/width
  \sbox\pandoc@box{#1}%
  \Gscale@div\@tempa{\textheight}{\dimexpr\ht\pandoc@box+\dp\pandoc@box\relax}%
  \Gscale@div\@tempb{\linewidth}{\wd\pandoc@box}%
  \ifdim\@tempb\p@<\@tempa\p@\let\@tempa\@tempb\fi% select the smaller of both
  \ifdim\@tempa\p@<\p@\scalebox{\@tempa}{\usebox\pandoc@box}%
  \else\usebox{\pandoc@box}%
  \fi%
}
% Set default figure placement to htbp
\def\fps@figure{htbp}
\makeatother



\ifLuaTeX
\usepackage[bidi=basic,shorthands=off,provide=*]{babel}
\else
\usepackage[bidi=default,shorthands=off,provide=*]{babel}
\fi
\ifPDFTeX
\else
\babelfont{rm}[]{Times New Roman}
\fi
\ifLuaTeX
  \usepackage{selnolig} % disable illegal ligatures
\fi


\setlength{\emergencystretch}{3em} % prevent overfull lines

\providecommand{\tightlist}{%
  \setlength{\itemsep}{0pt}\setlength{\parskip}{0pt}}



 


\makeatletter
\@ifpackageloaded{tcolorbox}{}{\usepackage[skins,breakable]{tcolorbox}}
\@ifpackageloaded{fontawesome5}{}{\usepackage{fontawesome5}}
\definecolor{quarto-callout-color}{HTML}{909090}
\definecolor{quarto-callout-note-color}{HTML}{0758E5}
\definecolor{quarto-callout-important-color}{HTML}{CC1914}
\definecolor{quarto-callout-warning-color}{HTML}{EB9113}
\definecolor{quarto-callout-tip-color}{HTML}{00A047}
\definecolor{quarto-callout-caution-color}{HTML}{FC5300}
\definecolor{quarto-callout-color-frame}{HTML}{acacac}
\definecolor{quarto-callout-note-color-frame}{HTML}{4582ec}
\definecolor{quarto-callout-important-color-frame}{HTML}{d9534f}
\definecolor{quarto-callout-warning-color-frame}{HTML}{f0ad4e}
\definecolor{quarto-callout-tip-color-frame}{HTML}{02b875}
\definecolor{quarto-callout-caution-color-frame}{HTML}{fd7e14}
\makeatother
\makeatletter
\@ifpackageloaded{caption}{}{\usepackage{caption}}
\AtBeginDocument{%
\ifdefined\contentsname
  \renewcommand*\contentsname{Содержание}
\else
  \newcommand\contentsname{Содержание}
\fi
\ifdefined\listfigurename
  \renewcommand*\listfigurename{Список Иллюстраций}
\else
  \newcommand\listfigurename{Список Иллюстраций}
\fi
\ifdefined\listtablename
  \renewcommand*\listtablename{Список Таблиц}
\else
  \newcommand\listtablename{Список Таблиц}
\fi
\ifdefined\figurename
  \renewcommand*\figurename{Рисунок}
\else
  \newcommand\figurename{Рисунок}
\fi
\ifdefined\tablename
  \renewcommand*\tablename{Таблица}
\else
  \newcommand\tablename{Таблица}
\fi
}
\@ifpackageloaded{float}{}{\usepackage{float}}
\floatstyle{ruled}
\@ifundefined{c@chapter}{\newfloat{codelisting}{h}{lop}}{\newfloat{codelisting}{h}{lop}[chapter]}
\floatname{codelisting}{Список}
\newcommand*\listoflistings{\listof{codelisting}{Список Каталогов}}
\makeatother
\makeatletter
\makeatother
\makeatletter
\@ifpackageloaded{caption}{}{\usepackage{caption}}
\@ifpackageloaded{subcaption}{}{\usepackage{subcaption}}
\makeatother
\usepackage{bookmark}
\IfFileExists{xurl.sty}{\usepackage{xurl}}{} % add URL line breaks if available
\urlstyle{same}
% fallback for those not using the hyperref driver hyperxmp:
\makeatletter
\@ifundefined{xmpquote}{\newcommand{\xmpquote}[1]{#1}}{}
\makeatother
\hypersetup{
  pdftitle={Задачи по Эконометрике-1: Степень подгонки линейной регрессии},
  pdfauthor={Н. В. Артамонов, А. Н. Лата (МГИМО МИД России)},
  pdflang={ru},
  colorlinks=true,
  linkcolor={blue},
  filecolor={Maroon},
  citecolor={Blue},
  urlcolor={Blue},
  pdfcreator={LaTeX via pandoc}}


\title{Задачи по Эконометрике-1: Степень подгонки линейной регрессии}
\author{Н. В. Артамонов, А. Н. Лата (МГИМО МИД России)}
\date{}
\begin{document}
\maketitle

\renewcommand*\contentsname{Содержание}
{
\setcounter{tocdepth}{3}
\tableofcontents
}

\begin{tcolorbox}[enhanced jigsaw, arc=.35mm, bottomrule=.15mm, bottomtitle=1mm, breakable, colback=white, colbacktitle=quarto-callout-important-color!10!white, colframe=quarto-callout-important-color-frame, coltitle=black, left=2mm, leftrule=.75mm, opacityback=0, opacitybacktitle=0.6, rightrule=.15mm, title=\textcolor{quarto-callout-important-color}{\faExclamation}\hspace{0.5em}{Важно}, titlerule=0mm, toprule=.15mm, toptitle=1mm]

\begin{itemize}
\tightlist
\item
  Везде log --- натуральный логарифм; константа включается в модель по
  умолчанию.
\item
  Наборы данных \texttt{sleep75}, \texttt{Labour}, \texttt{Electricity}
  лежат в папке \texttt{datasets/}, их описание --- в
  \texttt{datasets/data-description.md}.
\item
  Наблюдения нумеруются с нуля (\texttt{iloc}); в задачах на подвыборках
  --- порядковые номера внутри подвыборки.
\item
  Имена коэффициентов --- как в statsmodels: \texttt{Intercept} ---
  константа, \texttt{np.log(x)} --- \(\log x\), \texttt{I(x\ **\ 2)} ---
  \(x^2\).
\item
  Округление указано в каждом вопросе; ответы считаются по неокруглённым
  оценкам. Оценки, которые при округлении обращаются в ноль, приведены с
  тремя значащими цифрами. Расхождение с ответом на единицу последнего
  знака --- результат округления.
\end{itemize}

\end{tcolorbox}

\section{Суммы квадратов, остатки, предсказанные значения, коэффициент
детерминации}\label{ux441ux443ux43cux43cux44b-ux43aux432ux430ux434ux440ux430ux442ux43eux432-ux43eux441ux442ux430ux442ux43aux438-ux43fux440ux435ux434ux441ux43aux430ux437ux430ux43dux43dux44bux435-ux437ux43dux430ux447ux435ux43dux438ux44f-ux43aux43eux44dux444ux444ux438ux446ux438ux435ux43dux442-ux434ux435ux442ux435ux440ux43cux438ux43dux430ux446ux438ux438}

\subsection{sleep equation \#1}\label{sleep-equation-1}

Для набора данных \texttt{sleep75} рассмотрим регрессию \textbf{sleep на
totwrk, age, male, smsa}. Переменные \texttt{sleep} и \texttt{totwrk}
измерены в минутах в неделю.

Найдите коэффициенты подогнанной модели. \textbf{Ответ округлите до 2-х
десятичных знаков.}

{\def\LTcaptype{none} % do not increment counter
\begin{longtable}[]{@{}lr@{}}
\toprule\noalign{}
Переменная & Оценка \\
\midrule\noalign{}
\endhead
\bottomrule\noalign{}
\endlastfoot
Intercept & 3494.22 \\
totwrk & -0.17 \\
age & 2.81 \\
male & 86.91 \\
smsa & -75.29 \\
\end{longtable}
}

Модель была подогнана по 706 наблюдениям.

Для наблюдений с номерами 0, 3, 5, 9, 507, 699 (счёт с нуля,
\texttt{iloc}) найдите фактическое значение зависимой переменной,
предсказанное значение и остаток. \textbf{Ответ округлите до 3-х
десятичных знаков.}

{\def\LTcaptype{none} % do not increment counter
\begin{longtable}[]{@{}rrrr@{}}
\toprule\noalign{}
№ & Факт & Предсказание & Остаток \\
\midrule\noalign{}
\endhead
\bottomrule\noalign{}
\endlastfoot
0 & 3113 & 3094.454 & 18.546 \\
3 & 3083 & 2943.580 & 139.420 \\
5 & 4063 & 3494.143 & 568.857 \\
9 & 3018 & 3111.983 & -93.983 \\
507 & 3630 & 3476.914 & 153.086 \\
699 & 2993 & 3295.398 & -302.398 \\
\end{longtable}
}

Вычислите TSS, ESS, RSS. \textbf{Ответ округлите до 2-х десятичных
знаков.}

\begin{tcolorbox}[enhanced jigsaw, arc=.35mm, bottomrule=.15mm, bottomtitle=1mm, breakable, colback=white, colbacktitle=quarto-callout-note-color!10!white, colframe=quarto-callout-note-color-frame, coltitle=black, left=2mm, leftrule=.75mm, opacityback=0, opacitybacktitle=0.6, rightrule=.15mm, title=\textcolor{quarto-callout-note-color}{\faInfo}\hspace{0.5em}{Суммы квадратов}, titlerule=0mm, toprule=.15mm, toptitle=1mm]

\(TSS = 139239835.76\)

\(ESS = 17188382.23\)

\(RSS = 122051453.53\)

\end{tcolorbox}

Вычислите \(R^2\) и \(R^2_{adj}\). \textbf{Ответ округлите до 4-х
десятичных знаков.}

\begin{tcolorbox}[enhanced jigsaw, arc=.35mm, bottomrule=.15mm, bottomtitle=1mm, breakable, colback=white, colbacktitle=quarto-callout-note-color!10!white, colframe=quarto-callout-note-color-frame, coltitle=black, left=2mm, leftrule=.75mm, opacityback=0, opacitybacktitle=0.6, rightrule=.15mm, title=\textcolor{quarto-callout-note-color}{\faInfo}\hspace{0.5em}{Коэффициент детерминации}, titlerule=0mm, toprule=.15mm, toptitle=1mm]

\(R^2 = 0.1234\), \(R^2_{adj} = 0.1184\)

\end{tcolorbox}

Вычислите стандартную ошибку регрессии. \textbf{Ответ округлите до 2-х
десятичных знаков.}

\begin{tcolorbox}[enhanced jigsaw, arc=.35mm, bottomrule=.15mm, bottomtitle=1mm, breakable, colback=white, colbacktitle=quarto-callout-note-color!10!white, colframe=quarto-callout-note-color-frame, coltitle=black, left=2mm, leftrule=.75mm, opacityback=0, opacitybacktitle=0.6, rightrule=.15mm, title=\textcolor{quarto-callout-note-color}{\faInfo}\hspace{0.5em}{Стандартная ошибка регрессии}, titlerule=0mm, toprule=.15mm, toptitle=1mm]

417.27

\end{tcolorbox}

\subsection{sleep equation \#2}\label{sleep-equation-2}

Для набора данных \texttt{sleep75} рассмотрим регрессию \textbf{sleep на
totwrk, age, male, south} \emph{только для жителей мегаполиса (smsa =
1)}.

Найдите коэффициенты подогнанной модели. \textbf{Ответ округлите до 2-х
десятичных знаков.}

{\def\LTcaptype{none} % do not increment counter
\begin{longtable}[]{@{}lr@{}}
\toprule\noalign{}
Переменная & Оценка \\
\midrule\noalign{}
\endhead
\bottomrule\noalign{}
\endlastfoot
Intercept & 3460.43 \\
totwrk & -0.18 \\
age & 2.28 \\
male & 92.65 \\
south & 168.85 \\
\end{longtable}
}

Модель была подогнана по 282 наблюдениям.

Для наблюдений с порядковыми номерами внутри подвыборки 0, 8, 55, 75, 99
(счёт с нуля, \texttt{iloc}, после отбора строк с smsa = 1) найдите
фактическое значение зависимой переменной, предсказанное значение и
остаток. \textbf{Ответ округлите до 3-х десятичных знаков.}

{\def\LTcaptype{none} % do not increment counter
\begin{longtable}[]{@{}rrrrr@{}}
\toprule\noalign{}
№ & Индекс в sleep75 & Факт & Предсказание & Остаток \\
\midrule\noalign{}
\endhead
\bottomrule\noalign{}
\endlastfoot
0 & 2 & 2670 & 3132.491 & -462.491 \\
8 & 19 & 3090 & 3370.083 & -280.083 \\
55 & 149 & 2985 & 2835.360 & 149.640 \\
75 & 194 & 3660 & 3248.540 & 411.460 \\
99 & 259 & 2618 & 2863.070 & -245.070 \\
\end{longtable}
}

Вычислите TSS, ESS, RSS. \textbf{Ответ округлите до 2-х десятичных
знаков.}

\begin{tcolorbox}[enhanced jigsaw, arc=.35mm, bottomrule=.15mm, bottomtitle=1mm, breakable, colback=white, colbacktitle=quarto-callout-note-color!10!white, colframe=quarto-callout-note-color-frame, coltitle=black, left=2mm, leftrule=.75mm, opacityback=0, opacitybacktitle=0.6, rightrule=.15mm, title=\textcolor{quarto-callout-note-color}{\faInfo}\hspace{0.5em}{Суммы квадратов}, titlerule=0mm, toprule=.15mm, toptitle=1mm]

\(TSS = 50055043.40\)

\(ESS = 7396259.39\)

\(RSS = 42658784.01\)

\end{tcolorbox}

Вычислите \(R^2\) и \(R^2_{adj}\). \textbf{Ответ округлите до 4-х
десятичных знаков.}

\begin{tcolorbox}[enhanced jigsaw, arc=.35mm, bottomrule=.15mm, bottomtitle=1mm, breakable, colback=white, colbacktitle=quarto-callout-note-color!10!white, colframe=quarto-callout-note-color-frame, coltitle=black, left=2mm, leftrule=.75mm, opacityback=0, opacitybacktitle=0.6, rightrule=.15mm, title=\textcolor{quarto-callout-note-color}{\faInfo}\hspace{0.5em}{Коэффициент детерминации}, titlerule=0mm, toprule=.15mm, toptitle=1mm]

\(R^2 = 0.1478\), \(R^2_{adj} = 0.1355\)

\end{tcolorbox}

Вычислите стандартную ошибку регрессии. \textbf{Ответ округлите до 2-х
десятичных знаков.}

\begin{tcolorbox}[enhanced jigsaw, arc=.35mm, bottomrule=.15mm, bottomtitle=1mm, breakable, colback=white, colbacktitle=quarto-callout-note-color!10!white, colframe=quarto-callout-note-color-frame, coltitle=black, left=2mm, leftrule=.75mm, opacityback=0, opacitybacktitle=0.6, rightrule=.15mm, title=\textcolor{quarto-callout-note-color}{\faInfo}\hspace{0.5em}{Стандартная ошибка регрессии}, titlerule=0mm, toprule=.15mm, toptitle=1mm]

392.43

\end{tcolorbox}

\subsection{sleep equation \#3}\label{sleep-equation-3}

Для набора данных \texttt{sleep75} рассмотрим регрессию \textbf{sleep на
totwrk, age, smsa, south} \emph{только для мужчин (male = 1)}.

Найдите коэффициенты подогнанной модели. \textbf{Ответ округлите до 2-х
десятичных знаков.}

{\def\LTcaptype{none} % do not increment counter
\begin{longtable}[]{@{}lr@{}}
\toprule\noalign{}
Переменная & Оценка \\
\midrule\noalign{}
\endhead
\bottomrule\noalign{}
\endlastfoot
Intercept & 3600.96 \\
totwrk & -0.19 \\
age & 3.30 \\
smsa & -77.69 \\
south & 60.11 \\
\end{longtable}
}

Модель была подогнана по 400 наблюдениям.

Для наблюдений с порядковыми номерами внутри подвыборки 0, 8, 44, 186,
279 (счёт с нуля, \texttt{iloc}, после отбора строк с male = 1) найдите
фактическое значение зависимой переменной, предсказанное значение и
остаток. \textbf{Ответ округлите до 3-х десятичных знаков.}

{\def\LTcaptype{none} % do not increment counter
\begin{longtable}[]{@{}rrrrr@{}}
\toprule\noalign{}
№ & Индекс в sleep75 & Факт & Предсказание & Остаток \\
\midrule\noalign{}
\endhead
\bottomrule\noalign{}
\endlastfoot
0 & 0 & 3113 & 3061.481 & 51.519 \\
8 & 9 & 3018 & 3087.341 & -69.341 \\
44 & 48 & 3363 & 3268.521 & 94.479 \\
186 & 262 & 3390 & 3228.976 & 161.024 \\
279 & 440 & 2955 & 3399.060 & -444.060 \\
\end{longtable}
}

Вычислите TSS, ESS, RSS. \textbf{Ответ округлите до 2-х десятичных
знаков.}

\begin{tcolorbox}[enhanced jigsaw, arc=.35mm, bottomrule=.15mm, bottomtitle=1mm, breakable, colback=white, colbacktitle=quarto-callout-note-color!10!white, colframe=quarto-callout-note-color-frame, coltitle=black, left=2mm, leftrule=.75mm, opacityback=0, opacitybacktitle=0.6, rightrule=.15mm, title=\textcolor{quarto-callout-note-color}{\faInfo}\hspace{0.5em}{Суммы квадратов}, titlerule=0mm, toprule=.15mm, toptitle=1mm]

\(TSS = 75570140.58\)

\(ESS = 12125135.60\)

\(RSS = 63445004.98\)

\end{tcolorbox}

Вычислите \(R^2\) и \(R^2_{adj}\). \textbf{Ответ округлите до 4-х
десятичных знаков.}

\begin{tcolorbox}[enhanced jigsaw, arc=.35mm, bottomrule=.15mm, bottomtitle=1mm, breakable, colback=white, colbacktitle=quarto-callout-note-color!10!white, colframe=quarto-callout-note-color-frame, coltitle=black, left=2mm, leftrule=.75mm, opacityback=0, opacitybacktitle=0.6, rightrule=.15mm, title=\textcolor{quarto-callout-note-color}{\faInfo}\hspace{0.5em}{Коэффициент детерминации}, titlerule=0mm, toprule=.15mm, toptitle=1mm]

\(R^2 = 0.1604\), \(R^2_{adj} = 0.1519\)

\end{tcolorbox}

Вычислите стандартную ошибку регрессии. \textbf{Ответ округлите до 2-х
десятичных знаков.}

\begin{tcolorbox}[enhanced jigsaw, arc=.35mm, bottomrule=.15mm, bottomtitle=1mm, breakable, colback=white, colbacktitle=quarto-callout-note-color!10!white, colframe=quarto-callout-note-color-frame, coltitle=black, left=2mm, leftrule=.75mm, opacityback=0, opacitybacktitle=0.6, rightrule=.15mm, title=\textcolor{quarto-callout-note-color}{\faInfo}\hspace{0.5em}{Стандартная ошибка регрессии}, titlerule=0mm, toprule=.15mm, toptitle=1mm]

400.77

\end{tcolorbox}

\subsection{sleep equation \#4}\label{sleep-equation-4}

Для набора данных \texttt{sleep75} рассмотрим регрессию \textbf{sleep на
totwrk, age, smsa, south} \emph{только для женщин (male = 0)}.

Найдите коэффициенты подогнанной модели. \textbf{Ответ округлите до 2-х
десятичных знаков.}

{\def\LTcaptype{none} % do not increment counter
\begin{longtable}[]{@{}lr@{}}
\toprule\noalign{}
Переменная & Оценка \\
\midrule\noalign{}
\endhead
\bottomrule\noalign{}
\endlastfoot
Intercept & 3468.20 \\
totwrk & -0.15 \\
age & 1.65 \\
smsa & -37.69 \\
south & 136.03 \\
\end{longtable}
}

Модель была подогнана по 306 наблюдениям.

Для наблюдений с порядковыми номерами внутри подвыборки 0, 8, 44, 186,
279 (счёт с нуля, \texttt{iloc}, после отбора строк с male = 0) найдите
фактическое значение зависимой переменной, предсказанное значение и
остаток. \textbf{Ответ округлите до 3-х десятичных знаков.}

{\def\LTcaptype{none} % do not increment counter
\begin{longtable}[]{@{}rrrrr@{}}
\toprule\noalign{}
№ & Индекс в sleep75 & Факт & Предсказание & Остаток \\
\midrule\noalign{}
\endhead
\bottomrule\noalign{}
\endlastfoot
0 & 3 & 3083 & 2945.903 & 137.097 \\
8 & 67 & 2858 & 2920.662 & -62.662 \\
44 & 198 & 3730 & 3328.629 & 401.371 \\
186 & 472 & 3376 & 3503.722 & -127.722 \\
279 & 628 & 3405 & 3215.073 & 189.927 \\
\end{longtable}
}

Вычислите TSS, ESS, RSS. \textbf{Ответ округлите до 2-х десятичных
знаков.}

\begin{tcolorbox}[enhanced jigsaw, arc=.35mm, bottomrule=.15mm, bottomtitle=1mm, breakable, colback=white, colbacktitle=quarto-callout-note-color!10!white, colframe=quarto-callout-note-color-frame, coltitle=black, left=2mm, leftrule=.75mm, opacityback=0, opacitybacktitle=0.6, rightrule=.15mm, title=\textcolor{quarto-callout-note-color}{\faInfo}\hspace{0.5em}{Суммы квадратов}, titlerule=0mm, toprule=.15mm, toptitle=1mm]

\(TSS = 63490152.12\)

\(ESS = 6300990.57\)

\(RSS = 57189161.55\)

\end{tcolorbox}

Вычислите \(R^2\) и \(R^2_{adj}\). \textbf{Ответ округлите до 4-х
десятичных знаков.}

\begin{tcolorbox}[enhanced jigsaw, arc=.35mm, bottomrule=.15mm, bottomtitle=1mm, breakable, colback=white, colbacktitle=quarto-callout-note-color!10!white, colframe=quarto-callout-note-color-frame, coltitle=black, left=2mm, leftrule=.75mm, opacityback=0, opacitybacktitle=0.6, rightrule=.15mm, title=\textcolor{quarto-callout-note-color}{\faInfo}\hspace{0.5em}{Коэффициент детерминации}, titlerule=0mm, toprule=.15mm, toptitle=1mm]

\(R^2 = 0.0992\), \(R^2_{adj} = 0.0873\)

\end{tcolorbox}

Вычислите стандартную ошибку регрессии. \textbf{Ответ округлите до 2-х
десятичных знаков.}

\begin{tcolorbox}[enhanced jigsaw, arc=.35mm, bottomrule=.15mm, bottomtitle=1mm, breakable, colback=white, colbacktitle=quarto-callout-note-color!10!white, colframe=quarto-callout-note-color-frame, coltitle=black, left=2mm, leftrule=.75mm, opacityback=0, opacitybacktitle=0.6, rightrule=.15mm, title=\textcolor{quarto-callout-note-color}{\faInfo}\hspace{0.5em}{Стандартная ошибка регрессии}, titlerule=0mm, toprule=.15mm, toptitle=1mm]

435.89

\end{tcolorbox}

\subsection{output equation \#1}\label{output-equation-1}

Для набора данных \texttt{Labour} рассмотрим регрессию \textbf{output на
capital, labour, wage}. Переменные \texttt{output} и \texttt{capital}
измерены в млн евро, \texttt{wage} --- в тыс. евро на работника.

Найдите коэффициенты подогнанной модели. \textbf{Ответ округлите до 2-х
десятичных знаков.}

{\def\LTcaptype{none} % do not increment counter
\begin{longtable}[]{@{}lr@{}}
\toprule\noalign{}
Переменная & Оценка \\
\midrule\noalign{}
\endhead
\bottomrule\noalign{}
\endlastfoot
Intercept & -16.75 \\
capital & 0.40 \\
labour & 0.05 \\
wage & 0.44 \\
\end{longtable}
}

Модель была подогнана по 569 наблюдениям.

Для наблюдений с номерами 1, 10, 23, 100, 348, 500 (счёт с нуля,
\texttt{iloc}) найдите фактическое значение зависимой переменной,
предсказанное значение и остаток. \textbf{Ответ округлите до 3-х
десятичных знаков.}

{\def\LTcaptype{none} % do not increment counter
\begin{longtable}[]{@{}rrrr@{}}
\toprule\noalign{}
№ & Факт & Предсказание & Остаток \\
\midrule\noalign{}
\endhead
\bottomrule\noalign{}
\endlastfoot
1 & 3.664 & 0.457 & 3.207 \\
10 & 3.246 & 5.666 & -2.420 \\
23 & 2.666 & -0.625 & 3.292 \\
100 & 2.181 & 6.952 & -4.771 \\
348 & 4.307 & 2.278 & 2.029 \\
500 & 3.660 & 11.353 & -7.694 \\
\end{longtable}
}

Вычислите TSS, ESS, RSS. \textbf{Ответ округлите до 2-х десятичных
знаков.}

\begin{tcolorbox}[enhanced jigsaw, arc=.35mm, bottomrule=.15mm, bottomtitle=1mm, breakable, colback=white, colbacktitle=quarto-callout-note-color!10!white, colframe=quarto-callout-note-color-frame, coltitle=black, left=2mm, leftrule=.75mm, opacityback=0, opacitybacktitle=0.6, rightrule=.15mm, title=\textcolor{quarto-callout-note-color}{\faInfo}\hspace{0.5em}{Суммы квадратов}, titlerule=0mm, toprule=.15mm, toptitle=1mm]

\(TSS = 2235073.29\)

\(ESS = 2190376.08\)

\(RSS = 44697.22\)

\end{tcolorbox}

Вычислите \(R^2\) и \(R^2_{adj}\). \textbf{Ответ округлите до 4-х
десятичных знаков.}

\begin{tcolorbox}[enhanced jigsaw, arc=.35mm, bottomrule=.15mm, bottomtitle=1mm, breakable, colback=white, colbacktitle=quarto-callout-note-color!10!white, colframe=quarto-callout-note-color-frame, coltitle=black, left=2mm, leftrule=.75mm, opacityback=0, opacitybacktitle=0.6, rightrule=.15mm, title=\textcolor{quarto-callout-note-color}{\faInfo}\hspace{0.5em}{Коэффициент детерминации}, titlerule=0mm, toprule=.15mm, toptitle=1mm]

\(R^2 = 0.9800\), \(R^2_{adj} = 0.9799\)

\end{tcolorbox}

Вычислите стандартную ошибку регрессии. \textbf{Ответ округлите до 2-х
десятичных знаков.}

\begin{tcolorbox}[enhanced jigsaw, arc=.35mm, bottomrule=.15mm, bottomtitle=1mm, breakable, colback=white, colbacktitle=quarto-callout-note-color!10!white, colframe=quarto-callout-note-color-frame, coltitle=black, left=2mm, leftrule=.75mm, opacityback=0, opacitybacktitle=0.6, rightrule=.15mm, title=\textcolor{quarto-callout-note-color}{\faInfo}\hspace{0.5em}{Стандартная ошибка регрессии}, titlerule=0mm, toprule=.15mm, toptitle=1mm]

8.89

\end{tcolorbox}

\subsection{output equation \#2}\label{output-equation-2}

Для набора данных \texttt{Labour} рассмотрим регрессию
\textbf{np.log(output) на np.log(capital), np.log(labour),
np.log(wage)}.

Найдите коэффициенты подогнанной модели. \textbf{Ответ округлите до 2-х
десятичных знаков.}

{\def\LTcaptype{none} % do not increment counter
\begin{longtable}[]{@{}lr@{}}
\toprule\noalign{}
Переменная & Оценка \\
\midrule\noalign{}
\endhead
\bottomrule\noalign{}
\endlastfoot
Intercept & -5.01 \\
np.log(capital) & 0.15 \\
np.log(labour) & 0.72 \\
np.log(wage) & 0.92 \\
\end{longtable}
}

Модель была подогнана по 569 наблюдениям.

Для наблюдений с номерами 1, 10, 23, 100, 348, 500 (счёт с нуля,
\texttt{iloc}) найдите фактическое значение зависимой переменной (здесь
это \texttt{np.log(output)}), предсказанное значение и остаток.
\textbf{Ответ округлите до 3-х десятичных знаков.}

{\def\LTcaptype{none} % do not increment counter
\begin{longtable}[]{@{}rrrr@{}}
\toprule\noalign{}
№ & Факт & Предсказание & Остаток \\
\midrule\noalign{}
\endhead
\bottomrule\noalign{}
\endlastfoot
1 & 1.299 & 1.347 & -0.049 \\
10 & 1.177 & 1.427 & -0.250 \\
23 & 0.981 & 0.344 & 0.637 \\
100 & 0.780 & 1.006 & -0.226 \\
348 & 1.460 & 1.177 & 0.283 \\
500 & 1.297 & 1.730 & -0.433 \\
\end{longtable}
}

Вычислите TSS, ESS, RSS. \textbf{Ответ округлите до 2-х десятичных
знаков.}

\begin{tcolorbox}[enhanced jigsaw, arc=.35mm, bottomrule=.15mm, bottomtitle=1mm, breakable, colback=white, colbacktitle=quarto-callout-note-color!10!white, colframe=quarto-callout-note-color-frame, coltitle=black, left=2mm, leftrule=.75mm, opacityback=0, opacitybacktitle=0.6, rightrule=.15mm, title=\textcolor{quarto-callout-note-color}{\faInfo}\hspace{0.5em}{Суммы квадратов}, titlerule=0mm, toprule=.15mm, toptitle=1mm]

\(TSS = 797.67\)

\(ESS = 708.65\)

\(RSS = 89.02\)

\end{tcolorbox}

Вычислите \(R^2\) и \(R^2_{adj}\). \textbf{Ответ округлите до 4-х
десятичных знаков.}

\begin{tcolorbox}[enhanced jigsaw, arc=.35mm, bottomrule=.15mm, bottomtitle=1mm, breakable, colback=white, colbacktitle=quarto-callout-note-color!10!white, colframe=quarto-callout-note-color-frame, coltitle=black, left=2mm, leftrule=.75mm, opacityback=0, opacitybacktitle=0.6, rightrule=.15mm, title=\textcolor{quarto-callout-note-color}{\faInfo}\hspace{0.5em}{Коэффициент детерминации}, titlerule=0mm, toprule=.15mm, toptitle=1mm]

\(R^2 = 0.8884\), \(R^2_{adj} = 0.8878\)

\end{tcolorbox}

Вычислите стандартную ошибку регрессии. \textbf{Ответ округлите до 2-х
десятичных знаков.}

\begin{tcolorbox}[enhanced jigsaw, arc=.35mm, bottomrule=.15mm, bottomtitle=1mm, breakable, colback=white, colbacktitle=quarto-callout-note-color!10!white, colframe=quarto-callout-note-color-frame, coltitle=black, left=2mm, leftrule=.75mm, opacityback=0, opacitybacktitle=0.6, rightrule=.15mm, title=\textcolor{quarto-callout-note-color}{\faInfo}\hspace{0.5em}{Стандартная ошибка регрессии}, titlerule=0mm, toprule=.15mm, toptitle=1mm]

0.40

\end{tcolorbox}

\subsection{cost equation \#1}\label{cost-equation-1}

Для набора данных \texttt{Electricity} рассмотрим регрессию

\texttt{cost\ \textasciitilde{}\ q\ +\ I(q**2)\ +\ pk\ +\ pf\ +\ pl}

Найдите коэффициенты подогнанной модели. \textbf{Ответ округлите до 2-х
десятичных знаков.}

{\def\LTcaptype{none} % do not increment counter
\begin{longtable}[]{@{}lr@{}}
\toprule\noalign{}
Переменная & Оценка \\
\midrule\noalign{}
\endhead
\bottomrule\noalign{}
\endlastfoot
Intercept & -81.36 \\
q & 0.00432 \\
I(q ** 2) & 1.57e-08 \\
pk & 0.37 \\
pf & 0.95 \\
pl & 0.00361 \\
\end{longtable}
}

Модель была подогнана по 158 наблюдениям.

Для наблюдений с номерами 1, 16, 28, 47, 57, 93 (счёт с нуля,
\texttt{iloc}) найдите фактическое значение зависимой переменной,
предсказанное значение и остаток. \textbf{Ответ округлите до 3-х
десятичных знаков.}

{\def\LTcaptype{none} % do not increment counter
\begin{longtable}[]{@{}rrrr@{}}
\toprule\noalign{}
№ & Факт & Предсказание & Остаток \\
\midrule\noalign{}
\endhead
\bottomrule\noalign{}
\endlastfoot
1 & 3.043 & -2.512 & 5.555 \\
16 & 5.597 & 2.747 & 2.850 \\
28 & 7.439 & -2.269 & 9.709 \\
47 & 28.786 & 34.228 & -5.442 \\
57 & 41.902 & 43.941 & -2.040 \\
93 & 9.467 & -4.009 & 13.477 \\
\end{longtable}
}

Вычислите TSS, ESS, RSS. \textbf{Ответ округлите до 2-х десятичных
знаков.}

\begin{tcolorbox}[enhanced jigsaw, arc=.35mm, bottomrule=.15mm, bottomtitle=1mm, breakable, colback=white, colbacktitle=quarto-callout-note-color!10!white, colframe=quarto-callout-note-color-frame, coltitle=black, left=2mm, leftrule=.75mm, opacityback=0, opacitybacktitle=0.6, rightrule=.15mm, title=\textcolor{quarto-callout-note-color}{\faInfo}\hspace{0.5em}{Суммы квадратов}, titlerule=0mm, toprule=.15mm, toptitle=1mm]

\(TSS = 1189954.21\)

\(ESS = 1143993.96\)

\(RSS = 45960.25\)

\end{tcolorbox}

Вычислите \(R^2\) и \(R^2_{adj}\). \textbf{Ответ округлите до 4-х
десятичных знаков.}

\begin{tcolorbox}[enhanced jigsaw, arc=.35mm, bottomrule=.15mm, bottomtitle=1mm, breakable, colback=white, colbacktitle=quarto-callout-note-color!10!white, colframe=quarto-callout-note-color-frame, coltitle=black, left=2mm, leftrule=.75mm, opacityback=0, opacitybacktitle=0.6, rightrule=.15mm, title=\textcolor{quarto-callout-note-color}{\faInfo}\hspace{0.5em}{Коэффициент детерминации}, titlerule=0mm, toprule=.15mm, toptitle=1mm]

\(R^2 = 0.9614\), \(R^2_{adj} = 0.9601\)

\end{tcolorbox}

Вычислите стандартную ошибку регрессии. \textbf{Ответ округлите до 2-х
десятичных знаков.}

\begin{tcolorbox}[enhanced jigsaw, arc=.35mm, bottomrule=.15mm, bottomtitle=1mm, breakable, colback=white, colbacktitle=quarto-callout-note-color!10!white, colframe=quarto-callout-note-color-frame, coltitle=black, left=2mm, leftrule=.75mm, opacityback=0, opacitybacktitle=0.6, rightrule=.15mm, title=\textcolor{quarto-callout-note-color}{\faInfo}\hspace{0.5em}{Стандартная ошибка регрессии}, titlerule=0mm, toprule=.15mm, toptitle=1mm]

17.39

\end{tcolorbox}

\subsection{cost equation \#2}\label{cost-equation-2}

Для набора данных \texttt{Electricity} рассмотрим регрессию

\begin{verbatim}
np.log(cost) ~ np.log(q) + I(np.log(q)**2) +
    np.log(pk) + np.log(pf) + np.log(pl)
\end{verbatim}

Найдите коэффициенты подогнанной модели. \textbf{Ответ округлите до 2-х
десятичных знаков.}

{\def\LTcaptype{none} % do not increment counter
\begin{longtable}[]{@{}lr@{}}
\toprule\noalign{}
Переменная & Оценка \\
\midrule\noalign{}
\endhead
\bottomrule\noalign{}
\endlastfoot
Intercept & -6.74 \\
np.log(q) & 0.40 \\
I(np.log(q) ** 2) & 0.03 \\
np.log(pk) & 0.16 \\
np.log(pf) & 0.68 \\
np.log(pl) & 0.15 \\
\end{longtable}
}

Модель была подогнана по 158 наблюдениям.

Для наблюдений с номерами 1, 16, 28, 47, 57, 93 (счёт с нуля,
\texttt{iloc}) найдите фактическое значение зависимой переменной (здесь
это \texttt{np.log(cost)}), предсказанное значение и остаток.
\textbf{Ответ округлите до 3-х десятичных знаков.}

{\def\LTcaptype{none} % do not increment counter
\begin{longtable}[]{@{}rrrr@{}}
\toprule\noalign{}
№ & Факт & Предсказание & Остаток \\
\midrule\noalign{}
\endhead
\bottomrule\noalign{}
\endlastfoot
1 & 1.113 & 1.452 & -0.339 \\
16 & 1.722 & 1.650 & 0.072 \\
28 & 2.007 & 1.784 & 0.223 \\
47 & 3.360 & 3.556 & -0.196 \\
57 & 3.735 & 3.846 & -0.111 \\
93 & 2.248 & 2.247 & 0.001 \\
\end{longtable}
}

Вычислите TSS, ESS, RSS. \textbf{Ответ округлите до 2-х десятичных
знаков.}

\begin{tcolorbox}[enhanced jigsaw, arc=.35mm, bottomrule=.15mm, bottomtitle=1mm, breakable, colback=white, colbacktitle=quarto-callout-note-color!10!white, colframe=quarto-callout-note-color-frame, coltitle=black, left=2mm, leftrule=.75mm, opacityback=0, opacitybacktitle=0.6, rightrule=.15mm, title=\textcolor{quarto-callout-note-color}{\faInfo}\hspace{0.5em}{Суммы квадратов}, titlerule=0mm, toprule=.15mm, toptitle=1mm]

\(TSS = 373.66\)

\(ESS = 370.76\)

\(RSS = 2.90\)

\end{tcolorbox}

Вычислите \(R^2\) и \(R^2_{adj}\). \textbf{Ответ округлите до 4-х
десятичных знаков.}

\begin{tcolorbox}[enhanced jigsaw, arc=.35mm, bottomrule=.15mm, bottomtitle=1mm, breakable, colback=white, colbacktitle=quarto-callout-note-color!10!white, colframe=quarto-callout-note-color-frame, coltitle=black, left=2mm, leftrule=.75mm, opacityback=0, opacitybacktitle=0.6, rightrule=.15mm, title=\textcolor{quarto-callout-note-color}{\faInfo}\hspace{0.5em}{Коэффициент детерминации}, titlerule=0mm, toprule=.15mm, toptitle=1mm]

\(R^2 = 0.9922\), \(R^2_{adj} = 0.9920\)

\end{tcolorbox}

Вычислите стандартную ошибку регрессии. \textbf{Ответ округлите до 2-х
десятичных знаков.}

\begin{tcolorbox}[enhanced jigsaw, arc=.35mm, bottomrule=.15mm, bottomtitle=1mm, breakable, colback=white, colbacktitle=quarto-callout-note-color!10!white, colframe=quarto-callout-note-color-frame, coltitle=black, left=2mm, leftrule=.75mm, opacityback=0, opacitybacktitle=0.6, rightrule=.15mm, title=\textcolor{quarto-callout-note-color}{\faInfo}\hspace{0.5em}{Стандартная ошибка регрессии}, titlerule=0mm, toprule=.15mm, toptitle=1mm]

0.14

\end{tcolorbox}




\end{document}
